All 360 main certificates were accepted and replayed successfully at 192 bits. Every count agrees with its reference. The 36 delay counts agree with independent validated quadrature; all 12 single-delay cases also agree with the analytic crossing calculation below. The nine test methods pass, including explicit corrupted-record rejection and an unresolved boundary-zero example.

On the common-shift family, the raw certificate grows from 102 panels at $\kappa=0$ to 29,344 at $\kappa=512$. The anchor policy uses 102 panels, the midrange policy 84, and the local policy 28 for every shift. Thus the local policy's improvement from removing the common factor must be distinguished from its additional improvement over a valid global normalization. Figure~\ref{fig:shift} makes this distinction explicit.

Outside the five shift instances, local gauges produce fewer panels than Fixed-9 on 66 of 85 problems, tie on 19, and produce more panels on none. The median paired Fixed-9/local panel ratios are 2.17 for dominance problems and 1.85 for the generated lattice family. Their family median panel counts are shown in Table~\ref{tab:panels}. This is evidence about the tested finite grid and sufficient bound, not a proof that local gauges outperform every fixed global center on every function.

The practical delay families show a much smaller benefit against a well-chosen fixed center: their median paired panel ratios are about 1.01 and 1.02. Comparing only with midrange would overstate the practical effect, because the zero or anchor center is already effective on many delay problems. The clearance family is also nearly tied with Fixed-9: as $\epsilon$ decreases from $10^{-1}$ to $10^{-4}$, local counts rise from 36 to 76 panels. Gauge optimization cannot remove a zero approaching the boundary.

Table~\ref{tab:times} and Figure~\ref{fig:paired} separate construction from replay. In the lattice family, the median midrange/local construction ratio is 0.73, so local construction is slower, while the corresponding replay ratio is 1.97. For the single-delay family the ratios are 0.67 and 2.14, but these timing gains in replay are relative to the less effective midrange policy. No general runtime superiority follows. In particular, certificate length is the more portable outcome: it does not depend on system load, while timings do.
